% Naskah lengkap presentasi (materi utuh). Build: latexmk -pdf naskah.tex
\documentclass[11pt,a4paper]{article}
\usepackage[T1]{fontenc}
\usepackage[utf8]{inputenc}
\usepackage{lmodern}
\usepackage[indonesian]{babel}
\usepackage[margin=2.5cm]{geometry}
\usepackage{microtype}
\usepackage{amsmath,amssymb}
\usepackage{booktabs,tabularx,array,colortbl}
\usepackage{graphicx}
\usepackage{xcolor}
\usepackage{tikz}
\usetikzlibrary{arrows.meta,positioning,shapes.geometric,calc,fit,backgrounds}
\usepackage{pgfplots}
\usepackage{listings}
\usepackage[most]{tcolorbox}
\usepackage{caption}
\usepackage{enumitem}
\usepackage{titlesec}
\usepackage{needspace}
\usepackage{fancyhdr}
\usepackage[hidelinks]{hyperref}
% Shared colours, TikZ styles and pgfplots defaults (used by slides.tex and naskah.tex).
\definecolor{navy}{HTML}{1F3A5F}
\definecolor{accent}{HTML}{D9822B}
\definecolor{tealc}{HTML}{2A9D8F}
\definecolor{gapred}{HTML}{C8553D}
\definecolor{okgreen}{HTML}{3A8D5C}
\definecolor{midgray}{HTML}{8A94A6}
\definecolor{lightbg}{HTML}{EEF2F7}
\definecolor{insuf}{HTML}{C9CED8}

\pgfplotsset{
  compat=1.18,
  every axis/.append style={
    font=\scriptsize,
    axis line style={midgray},
    tick style={midgray},
    label style={font=\scriptsize},
    tick label style={font=\scriptsize},
    legend style={font=\scriptsize, draw=none, fill=none},
    grid style={midgray!25},
    every node near coord/.append style={text=black!80},
  },
}
% Indonesian number format everywhere (0,5 and 1.000)
\pgfkeys{/pgf/number format/.cd, use comma, 1000 sep={.}}

% ragged-right columns (avoid stretched inter-word spacing in narrow cells)
\newcolumntype{P}[1]{>{\raggedright\arraybackslash}p{#1}}
\newcolumntype{L}{>{\raggedright\arraybackslash}X}

\tikzset{
  box/.style={draw=navy, rounded corners=3pt, fill=white, align=center, font=\scriptsize,
              inner sep=4pt, line width=0.6pt},
  boxnavy/.style={box, fill=navy, text=white},
  boxacc/.style={box, draw=accent, fill=accent, text=white},
  boxteal/.style={box, draw=tealc, fill=tealc, text=white},
  boxsoft/.style={box, draw=midgray, fill=lightbg},
  arr/.style={-{Stealth[length=2.2mm]}, line width=0.8pt, draw=navy!80},
  lbl/.style={font=\tiny, text=black!70, align=center},
}

% height used by every pgfplots figure; set before \input{figs/...}
\newlength\figH
\setlength\figH{5cm}


\captionsetup{font=small, labelfont={bf,color=navy}}
% keep headings with their first lines (no heading alone at the bottom of a page)
\titleformat{\section}{\Needspace{10\baselineskip}\large\bfseries\color{navy}}{\thesection}{0.6em}{}
\titleformat{\subsection}{\Needspace{7\baselineskip}\normalsize\bfseries\color{navy}}{\thesubsection}{0.6em}{}
\setlist{itemsep=2pt, topsep=4pt}
\pagestyle{fancy}
\fancyhf{}
\fancyhead[L]{\small\color{midgray}Naskah presentasi -- P-DINA + Gemini LLM}
\fancyhead[R]{\small\color{midgray}\thepage}
\renewcommand{\headrulewidth}{0.3pt}
\setlength{\parskip}{4pt}

\newtcolorbox{catatan}[1][]{colback=lightbg, colframe=navy, boxrule=0.4pt, arc=2pt, left=6pt, right=6pt,
  top=4pt, bottom=4pt, fontupper=\small, #1}
\newtcolorbox{peringatan}[1][]{colback=accent!8, colframe=accent, boxrule=0.4pt, arc=2pt, left=6pt, right=6pt,
  top=4pt, bottom=4pt, fontupper=\small, #1}
\newcommand{\slideref}[1]{\marginpar{\raggedright\tiny\color{midgray}#1}}
\newcommand{\sumber}[1]{{\footnotesize\color{midgray}Sumber: #1}}
\lstset{basicstyle=\ttfamily\scriptsize, breaklines=true, columns=fullflexible, keepspaces=true,
  backgroundcolor=\color{lightbg}, frame=none, xleftmargin=4pt, xrightmargin=4pt}

\title{\color{navy}\bfseries Diagnosis Learning Gap Matematika dengan P-DINA\\ dan Laporan Naratif LLM\\[4pt]
  \large\mdseries Naskah lengkap presentasi}
\author{Muhammad Fariq Faqih}
\date{\today}

\begin{document}
\maketitle
\thispagestyle{empty}

\begin{catatan}
  \textbf{Cara memakai naskah ini.} Setiap bagian mengikuti satu slide (atau satu kelompok slide) pada
  \texttt{slides.pdf}, dengan urutan yang sama. Teks dapat dibacakan langsung atau diringkas. Semua angka berasal dari
  artefak hasil run proyek (\texttt{artifacts/}) dan dari satu laporan Gemini yang dibuat pada 27 Oktober 2026; tidak
  ada angka perkiraan.
\end{catatan}

\begin{peringatan}
  \textbf{Status data.} EEDI NeurIPS 2020 Task 3/4 dipakai sebagai \emph{data pengembangan dan evaluasi awal}.
  Populasinya mayoritas siswa SMP di Inggris, dan konsep targetnya hanya prasyarat aljabar SMA tanpa kalkulus. Hasil di
  naskah ini tidak boleh disajikan sebagai validasi pada kelas SMA penelitian.
\end{peringatan}

\tableofcontents
\newpage

% ====================================================================================================
\section{Pembuka dan agenda}

Selamat pagi/siang. Presentasi ini membahas sistem untuk \emph{mendiagnosis learning gap} siswa pada materi
matematika dan memberi rekomendasi intervensi. Sistem terdiri atas dua model dengan peran berbeda: \textbf{P-DINA},
model diagnosis kognitif yang menghitung probabilitas penguasaan konsep, dan \textbf{Gemini 3.1 Flash-Lite}, model
bahasa besar yang hanya bertugas menyampaikan hasil diagnosis itu dalam bahasa Indonesia yang mudah dibaca guru dan
siswa.

Urutan pembahasan: posisi kedua model di pipeline, pipeline end-to-end, konsep P-DINA, parameter awal dan proses
training, data dan evaluasi performa beserta alasan memilih P-DINA, hyperparameter tuning, penyimpanan model,
pemilihan dan spesifikasi LLM, konfigurasinya, system prompt, contoh output sistem secara keseluruhan, lalu
kesimpulan, keterbatasan, dan langkah berikutnya.

% ====================================================================================================
\section{Posisi tahap model di pipeline}

Hal pertama yang perlu jelas adalah \emph{siapa mengerjakan apa}. Gambar~\ref{fig:posisi} memperlihatkan dua tahap.
Tahap \textbf{komputasi} (biru) bersifat deterministik dan dapat diaudit: respons siswa masuk ke P-DINA, lalu P-DINA
menghasilkan probabilitas penguasaan, label, learning gap, dan rekomendasi berbasis aturan. Tahap \textbf{komunikasi}
(oranye) bersifat opsional: hasil diagnosis diringkas menjadi fakta JSON, dikirim ke Gemini untuk disusun menjadi
kalimat, lalu diperiksa validator sebelum ditampilkan.

\begin{figure}[h]
  \centering
  \resizebox{\linewidth}{!}{% Where the two models sit: P-DINA computes, Gemini only communicates.
\begin{tikzpicture}[node distance=5mm and 6mm]
  \node[boxsoft, minimum width=2.1cm, minimum height=1.1cm] (resp) {Respons siswa\\[1pt]{\tiny benar / salah per soal}};
  \node[boxnavy, right=of resp, minimum width=2.4cm, minimum height=1.1cm] (pdina) {\textbf{P-DINA}\\[1pt]{\tiny menghitung P(mastery)}};
  \node[boxsoft, right=of pdina, minimum width=2.5cm, minimum height=1.1cm] (diag) {Diagnosis\\[1pt]{\tiny label, learning gap,}\\[-1pt]{\tiny rekomendasi R1--R5}};
  \node[boxsoft, right=of diag, minimum width=2.3cm, minimum height=1.1cm] (facts) {Fakta JSON\\[1pt]{\tiny agregat, ID pseudonim}};
  \node[boxacc, right=of facts, minimum width=2.4cm, minimum height=1.1cm] (llm) {\textbf{Gemini LLM}\\[1pt]{\tiny menyusun kalimat}};
  \node[boxsoft, below=8mm of llm, minimum width=2.4cm] (val) {Validator\\[1pt]{\tiny angka, heading, klaim}};
  \node[boxsoft, below=8mm of diag, minimum width=2.5cm] (app) {Aplikasi Streamlit\\[1pt]{\tiny tabel, grafik, laporan}};

  \draw[arr] (resp) -- (pdina);
  \draw[arr] (pdina) -- (diag);
  \draw[arr] (diag) -- (facts);
  \draw[arr] (facts) -- (llm);
  \draw[arr] (llm) -- (val);
  \draw[arr] (val) -- node[lbl, above]{lolos / template} (app);
  \draw[arr] (diag) -- node[lbl, right]{langsung} (app);

  \begin{scope}[on background layer]
    \node[fill=navy!7, rounded corners, fit=(resp)(pdina)(diag), inner sep=6pt,
          label={[font=\tiny\bfseries, text=navy]above:Tahap komputasi (deterministik, dapat diaudit)}] {};
    \node[fill=accent!10, rounded corners, fit=(facts)(llm)(val), inner sep=6pt,
          label={[font=\tiny\bfseries, text=accent!80!black]above:Tahap komunikasi (opsional)}] {};
  \end{scope}
\end{tikzpicture}
}
  \caption{Posisi P-DINA dan Gemini di pipeline. Panah ``langsung'' menunjukkan bahwa tabel dan grafik di aplikasi
  tidak bergantung pada LLM.}
  \label{fig:posisi}
\end{figure}

\begin{table}[h]
  \centering\small
  \begin{tabularx}{\linewidth}{@{}l L L@{}}
    \toprule
    & \textbf{P-DINA} & \textbf{Gemini LLM} \\
    \midrule
    Peran & menghitung P(mastery) per siswa $\times$ konsep & menyusun laporan berbahasa Indonesia \\
    Input & respons benar/salah dan Q-matrix & fakta JSON ringkas hasil P-DINA, ID siswa dipseudonimkan \\
    Wajib? & ya, inti sistem & tidak; tanpa API key sistem menampilkan laporan template \\
    Dilarang & -- & menghitung atau mengubah angka, menambah konsep, mengklaim intervensi efektif \\
    \bottomrule
  \end{tabularx}
  \caption{Pembagian peran kedua model.}
\end{table}

Pemisahan ini disengaja. Keputusan pendidikan (siapa butuh remediasi, konsep mana yang diprioritaskan) harus berasal
dari model yang dapat diperiksa langkahnya. LLM tidak pernah menulis balik ke hasil diagnosis; bila LLM gagal atau
tidak tersedia, seluruh fungsi diagnosis tetap berjalan.

% ====================================================================================================
\section{Pipeline end-to-end}

Gambar~\ref{fig:pipeline} merangkum enam tahap sistem.

\begin{figure}[h]
  \centering
  \resizebox{\linewidth}{!}{% End-to-end pipeline: six stages.
\begin{tikzpicture}[
  stage/.style={box, minimum width=2.15cm, minimum height=2.5cm, text width=2.0cm, anchor=north},
  num/.style={circle, fill=#1, text=white, font=\scriptsize\bfseries, inner sep=1.5pt, minimum size=5mm},
]
  \foreach \xpos/\col/\ttl/\desc [count=\n] in {
      0/navy/{Data EEDI}/{1.382.727 jawaban\\4.918 siswa\\948 soal, 57 konsep},
      2.45/navy/{Audit \& Q-matrix}/{skema, leakage,\\soal\,$\to$\,konsep,\\11 konsep target},
      4.9/navy/{Training P-DINA}/{variational EM\\split per siswa\\80/10/10},
      7.35/navy/{Diagnosis \& rekomendasi}/{P(mastery), label,\\learning gap,\\aturan R1--R5},
      9.8/accent/{Laporan naratif LLM}/{Gemini 3.1\\Flash-Lite +\\validator},
      12.25/tealc/{Aplikasi Streamlit}/{asesmen, gap,\\intervensi,\\pre/post}} {
    \node[stage, draw=\col] (s\n) at (\xpos,0) {\textbf{\ttl}\\[3pt]{\tiny\desc}};
    \node[num=\col] at (s\n.north) {\n};
  }
  \foreach \a/\b in {1/2,2/3,3/4,4/5,5/6} {
    \draw[arr] ($(s\a.east)+(0,-1.25)$) -- ($(s\b.west)+(0,-1.25)$);
  }
\end{tikzpicture}
}
  \caption{Enam tahap pipeline. Tahap 1--4 menghasilkan checkpoint dan diagnosis; tahap 5 hanya membaca hasilnya.}
  \label{fig:pipeline}
\end{figure}

\begin{enumerate}
  \item \textbf{Data EEDI.} Arsip resmi NeurIPS 2020 Task 3/4 berisi 1.382.727 jawaban dari 4.918 siswa pada 948
        soal pilihan ganda. Hasil preprocessing dipakai \emph{penuh}: setelah pembersihan (opsi jawaban valid,
        satu jawaban terakhir per pasangan siswa--soal) tidak ada baris yang terbuang dan tidak ada sampling.
  \item \textbf{Audit dan Q-matrix.} Skema data diperiksa, kolom yang berisiko membocorkan jawaban dikecualikan, dan
        soal dipetakan ke 57 konsep (Q-matrix 948 $\times$ 57). Sebanyak 11 konsep aljabar dipilih sebagai target.
  \item \textbf{Training P-DINA.} Siswa dibagi 80/10/10 menjadi train, validation, dan test. P-DINA dilatih dengan
        variational EM pada data train.
  \item \textbf{Diagnosis dan rekomendasi.} Untuk setiap siswa dan konsep dihitung P(mastery), label, learning gap,
        prioritas, dan rekomendasi intervensi berbasis aturan transparan R1--R5.
  \item \textbf{Laporan naratif LLM.} Fakta diagnosis disusun menjadi laporan guru atau siswa oleh Gemini 3.1
        Flash-Lite, lalu divalidasi.
  \item \textbf{Aplikasi Streamlit.} Guru memakai semua fungsi lewat aplikasi: asesmen, analisis gap, rekomendasi
        intervensi, dan analisis pre-test/post-test.
\end{enumerate}

Prinsip penting di sepanjang pipeline: data train dipakai untuk fit, validation untuk seleksi, dan test hanya untuk
evaluasi \emph{satu kali}. Aplikasi tidak pernah melatih ulang; aplikasi hanya memuat checkpoint.

% ====================================================================================================
\section{P-DINA: probabilistic DINA dengan variational EM}

\subsection{Jenis masalah}

Masalah yang diselesaikan adalah \emph{diagnosis kognitif}: dari pola jawaban benar/salah, perkirakan probabilitas
seorang siswa menguasai setiap konsep. Ada tiga ciri yang membuat masalah ini berbeda dari klasifikasi biasa.

\begin{itemize}
  \item \textbf{Variabel laten.} Penguasaan $\alpha_{ik}\in\{0,1\}$ tidak pernah diamati. Tidak ada label ``siswa ini
        sebenarnya menguasai konsep ini'' yang bisa dipakai sebagai target supervised learning.
  \item \textbf{Multi-atribut.} Ke-57 konsep diperkirakan sekaligus dan saling terhubung melalui Q-matrix, yaitu tabel
        biner dari pakar yang menyatakan soal mana menguji konsep mana.
  \item \textbf{Data jarang.} Setiap siswa hanya mengerjakan sebagian kecil soal. Soal yang tidak dikerjakan tidak
        boleh dianggap salah.
\end{itemize}

\subsection{Model formal}

Intuisi DINA (\emph{Deterministic Inputs, Noisy ``And'' gate}) sederhana: siswa yang menguasai semua konsep yang
dibutuhkan sebuah soal hampir selalu menjawab benar, kecuali ceroboh (\emph{slip}, peluang $s_j$); siswa yang belum
menguasai biasanya salah, kecuali menebak benar (\emph{guess}, peluang $g_j$). Untuk siswa $i$, soal $j$, dan konsep
$k$:
\begin{align}
  \eta_{ij} &= \prod_k \alpha_{ik}^{\,q_{jk}}, \\
  P(X_{ij}=1\mid\boldsymbol\alpha_i) &= (1-s_j)^{\eta_{ij}}\, g_j^{\,1-\eta_{ij}}, \\
  \alpha_{ik} &\sim \mathrm{Bernoulli}\big(\sigma(a_k\theta_i-b_k)\big),\qquad \theta_i\sim\mathcal N(0,1).
\end{align}
Persamaan (3) adalah prior \emph{higher-order}: kemampuan umum $\theta_i$ membuat siswa yang kuat di banyak konsep
cenderung kuat di konsep lain, sehingga konsep dengan sedikit bukti dapat ``meminjam'' informasi dari konsep lain
milik siswa yang sama.

EM DINA klasik harus menjumlahkan semua profil penguasaan, yaitu $2^{57}\approx1{,}4\times10^{17}$ kemungkinan per
siswa, mustahil dihitung. P-DINA memakai pendekatan variasional \emph{mean-field}
$q(\boldsymbol\alpha_i)=\prod_k\mathrm{Bernoulli}(\pi_{ik})$, sehingga batas bawah log-likelihood (ELBO) punya bentuk
tertutup dan dapat dioptimalkan dengan perkalian matriks.

\subsection{Keunggulan}

\begin{table}[h]
  \centering\small
  \begin{tabularx}{\linewidth}{@{}>{\bfseries}l L@{}}
    \toprule
    Keunggulan & Penjelasan \\
    \midrule
    Interpretabel & setiap soal punya $g$ dan $s$; setiap konsep punya P(mastery) yang langsung bermakna bagi guru \\
    Ringan & sekitar 2.000 parameter; satu fit sekitar 30 detik di CPU laptop \\
    Missing $\neq$ salah & likelihood hanya menjumlahkan respons yang teramati \\
    Tahu kapan tidak tahu & bila bukti sedikit, label yang keluar \emph{insufficient\_evidence}, bukan learning gap \\
    Terkalibrasi & ECE test 0,008, dibanding 0,036 pada baseline Q-logistic \\
    Siswa baru & cukup E-step dengan parameter beku; diagnosis 30 siswa sekitar 0,4 detik \\
    Update eksak & untuk 934 dari 948 soal (soal satu-konsep), E-step sama dengan posterior Bayes tepat \\
    \bottomrule
  \end{tabularx}
  \caption{Keunggulan P-DINA untuk konteks pendidikan.}
\end{table}

Satu catatan jujur: kualitas diagnosis bergantung pada ketepatan Q-matrix. Audit visual terhadap 22 soal fokus
menunjukkan 17 soal cocok dengan label konsepnya dan 5 hanya cocok sebagian. Kesalahan Q-matrix langsung menjadi
kesalahan diagnosis.

% ====================================================================================================
\section{Parameter awal dan proses training}

Gambar~\ref{fig:em} memperlihatkan alur training. Semua parameter soal diinisialisasi pada nilai prior yang netral,
lalu E-step dan M-step dijalankan bergantian sampai ELBO berhenti naik.

\begin{figure}[h]
  \centering
  \resizebox{0.95\linewidth}{!}{% Variational EM training loop.
\begin{tikzpicture}[node distance=4mm and 5mm]
  \node[boxsoft, text width=2.3cm] (init) {\textbf{Inisialisasi}\\{\tiny $g_j{=}s_j{=}0{,}2$, $a_k{=}1$, $b_k{=}0$, $\theta_i{=}0$}};
  \node[boxnavy, right=of init, text width=2.5cm] (e) {\textbf{E-step}\\{\tiny perbarui $\pi_{ik}$ (keyakinan mastery) dan $\theta_i$ (Newton)}};
  \node[boxnavy, right=of e, text width=2.5cm] (m) {\textbf{M-step}\\{\tiny $g_j, s_j$ tertutup (MAP Beta); $a_k, b_k$ (Newton)}};
  \node[boxsoft, right=of m, text width=2.1cm] (elbo) {\textbf{Hitung ELBO}};
  \node[diamond, draw=navy, aspect=2.2, inner sep=1pt, font=\tiny, align=center, right=of elbo] (conv) {$\Delta\text{ELBO}/|\text{ELBO}|$\\$\le 10^{-6}$?};
  \node[boxacc, below=7mm of conv, text width=1.9cm] (stop) {\textbf{Checkpoint}\\{\tiny $g, s, a, b$}};
  \draw[arr] (init) -- (e);
  \draw[arr] (e) -- (m);
  \draw[arr] (m) -- (elbo);
  \draw[arr] (elbo) -- (conv);
  \draw[arr] (conv) -- node[lbl, right]{ya (atau 150 iterasi)} (stop);
  \draw[arr] (conv.north) -- ++(0,0.55) -| node[lbl, pos=0.25, above]{tidak} (e.north);
\end{tikzpicture}
}
  \caption{Loop variational EM.}
  \label{fig:em}
\end{figure}

\begin{table}[h]
  \centering\small
  \begin{tabularx}{\linewidth}{@{}l l L@{}}
    \toprule
    Parameter & Nilai awal & Keterangan \\
    \midrule
    guess $g_j$, slip $s_j$ (948 soal) & 0,2 & prior $\mathrm{Beta}(1+0{,}2\kappa;\ 1+0{,}8\kappa)$; terpilih $\kappa=2$, yaitu $\mathrm{Beta}(1{,}4;\ 2{,}6)$ \\
    $a_k$, $b_k$ (57 konsep) & 1 dan 0 & parameter prior higher-order \\
    $\theta_i$ (3.934 siswa train) & 0 & kemampuan umum, estimasi MAP dengan Newton \\
    $\pi_{ik}$ & dari E-step pertama & keyakinan bahwa siswa $i$ menguasai konsep $k$ \\
    batas & $g,s\in[0{,}001;\ 0{,}7]$ & dan $g+s\le0{,}98$, agar soal lebih mudah bagi yang menguasai \\
    \bottomrule
  \end{tabularx}
  \caption{Parameter awal P-DINA (\texttt{cdacdm/models.py}).}
\end{table}

\begin{itemize}
  \item \textbf{E-step} memperbarui $\pi_{ik}$. Untuk soal satu-konsep, update bersifat eksak dan dihitung dengan dua
        perkalian matriks. Untuk 14 soal multi-konsep dipakai iterasi titik tetap dengan \emph{damping} 0,5. Pada prior
        higher-order, $\theta_i$ diperbarui dengan metode Newton.
  \item \textbf{M-step} memperbarui $g_j$ dan $s_j$ dengan rumus tertutup (MAP dengan prior Beta), lalu $a_k$ dan $b_k$
        dengan Newton.
  \item \textbf{Kriteria berhenti}: perubahan relatif ELBO $\le10^{-6}$ atau 150 iterasi.
\end{itemize}

Model terpilih konvergen dalam \textbf{44 iterasi EM} dengan ELBO akhir $-630.546$, dalam \textbf{29,8 detik} di CPU
laptop.

\subsection{Parameter yang dipelajari}

\begin{figure}[h]
  \centering
  \setlength\figH{5.2cm}
  \begin{minipage}{0.62\linewidth}% Learned guess/slip distribution for 948 items; artifacts/model/item_parameters.csv (bins of 0.05)
\begin{tikzpicture}
\begin{axis}[
  ybar, width=\linewidth, height=\figH, bar width=3.2pt,
  xmin=-0.02, xmax=0.72, ymin=0,
  xlabel={nilai parameter}, ylabel={jumlah soal},
  ymajorgrids, legend style={at={(0.5,1.03)}, anchor=south, legend columns=2}, legend cell align=left,
  xticklabel style={/pgf/number format/fixed, /pgf/number format/precision=1},
]
  \addplot+[fill=navy, draw=none] coordinates {
    (0.025,3) (0.075,9) (0.125,38) (0.175,67) (0.225,129) (0.275,165) (0.325,162)
    (0.375,106) (0.425,81) (0.475,74) (0.525,46) (0.575,24) (0.625,21) (0.675,23)};
  \addlegendentry{guess $g$ (rata-rata 0,338)}
  \addplot+[fill=accent, draw=none] coordinates {
    (0.025,101) (0.075,141) (0.125,125) (0.175,86) (0.225,74) (0.275,79) (0.325,71)
    (0.375,57) (0.425,57) (0.475,46) (0.525,36) (0.575,25) (0.625,19) (0.675,31)};
  \addlegendentry{slip $s$ (rata-rata 0,253)}
\end{axis}
\end{tikzpicture}
\end{minipage}\hfill
  \begin{minipage}{0.35\linewidth}\small
    \begin{tabular}{@{}lccc@{}}
      \toprule
      & rata-rata & median \\
      \midrule
      $g$ & 0,338 & 0,317 \\
      $s$ & 0,253 & 0,217 \\
      $1-s-g$ & 0,409 & 0,417 \\
      \bottomrule
    \end{tabular}
  \end{minipage}
  \caption{Sebaran guess dan slip untuk 948 soal (\texttt{artifacts/model/item\_parameters.csv}).}
\end{figure}

Rata-rata guess 0,338 wajar untuk pilihan ganda empat opsi (tebakan acak 0,25). Daya beda $1-s-g$ rata-rata 0,409
tergolong sedang, sehingga satu konsep membutuhkan sekitar lima respons agar diagnosisnya stabil. Sebanyak 900 dari
948 soal punya parameter andal (minimal 30 respons train). Rata-rata $a_k=3{,}67$ menunjukkan penguasaan antar-konsep
sangat terkait dengan kemampuan umum.

% ====================================================================================================
\section{Data, evaluasi performa, dan alasan menggunakan P-DINA}

\subsection{Data dan protokol}

\begin{table}[h]
  \centering\small
  \begin{tabular}{@{}lrrl@{}}
    \toprule
    Split & Siswa & Interaksi & Dipakai untuk \\
    \midrule
    train & 3.934 & 1.112.450 & fit P-DINA dan baseline \\
    validation & 492 & 129.242 & seleksi hyperparameter dan kebijakan diagnosis \\
    test & 492 & 141.035 & evaluasi akhir, satu kali \\
    \midrule
    total & 4.918 & 1.382.727 & \\
    \bottomrule
  \end{tabular}
  \caption{Pembagian data. Split dilakukan per siswa (seed 42), sehingga tidak ada siswa di dua split.}
\end{table}

Artefak data preparation dibangun ulang secara lokal dari arsip resmi, karena notebook preprocessing asli memanggil
fungsi yang tidak didefinisikan dan file Drive tidak dapat diakses dari lingkungan pengembangan. Hasil rekonstruksi
dibandingkan dengan salinan file Drive: pada train dan validation, siswanya sama, tidak ada pasangan siswa--soal yang
hilang, dan tidak ada satu pun nilai benar/salah yang berbeda.

Model hanya membaca tiga kolom: \texttt{UserId}, \texttt{QuestionId}, dan \texttt{IsCorrect}. Kolom seperti pass rate
soal, difficulty, dan prior accuracy dikecualikan karena dapat membocorkan jawaban.

Protokol evaluasinya \emph{held-out}: pada setiap siswa validation atau test, 20\% respons disembunyikan (minimal lima
respons tetap terlihat sebagai konteks). Parameter soal dibekukan dari train, penguasaan ditebak dari konteks, lalu
respons yang disembunyikan diprediksi. Interval kepercayaan 95\% dihitung dengan bootstrap per siswa sebanyak 200
ulangan.

\begin{peringatan}
  EEDI tidak memiliki ground truth mastery. Prediksi respons yang baik tidak membuktikan diagnosis benar, sehingga
  evaluasi dilengkapi konsistensi split-half dan simulasi.
\end{peringatan}

\subsection{Hasil prediksi respons pada test}

\begin{figure}[h]
  \centering
  \setlength\figH{4.2cm}
  % Test ROC-AUC with bootstrap 95% CI (left) and ECE (right); artifacts/evaluation/test_metrics.csv
% Rows (bottom -> top): 1 item-mean, 2 concept-score, 3 Q-logistic, 4 P-DINA independent, 5 P-DINA main
\begin{tikzpicture}
\pgfplotsset{
  rowbars/.style={xbar, bar width=7pt, ymin=0.4, ymax=5.6, ytick={1,...,5}, xmajorgrids, height=\figH,
                  scaled x ticks=false},
  ci/.style={draw=none, bar shift=0pt, fill=#1,
             error bars/.cd, x dir=both, x explicit, error bar style={black!70}},
}
\begin{axis}[
  name=auc, rowbars, scale only axis, width=0.42\linewidth,
  xmin=0.62, xmax=0.80,
  yticklabels={Item-mean, Concept-score, Q-logistic, P-DINA independen, \textbf{P-DINA (utama)}},
  xlabel={ROC-AUC test [CI 95\%]},
  xticklabel style={/pgf/number format/fixed, /pgf/number format/precision=2},
]
  \addplot+[ci=midgray]  coordinates {(0.6722,1) += (0.0095,0) -= (0.0114,0)};
  \addplot+[ci=midgray]  coordinates {(0.7182,2) += (0.0148,0) -= (0.0114,0)};
  \addplot+[ci=tealc]    coordinates {(0.7746,3) += (0.0116,0) -= (0.0089,0)};
  \addplot+[ci=navy!45]  coordinates {(0.7613,4) += (0.0102,0) -= (0.0099,0)};
  \addplot+[ci=navy]     coordinates {(0.7684,5) += (0.0110,0) -= (0.0097,0)};
\end{axis}
\begin{axis}[
  at={(auc.east)}, anchor=west, xshift=0.5cm, rowbars, scale only axis, width=0.27\linewidth,
  xmin=0, xmax=0.045, yticklabels={},
  xlabel={ECE (makin kecil makin baik)},
  xtick={0,0.01,...,0.04},
  xticklabel style={/pgf/number format/fixed, /pgf/number format/precision=2},
  nodes near coords, point meta=explicit symbolic,
]
  \addplot+[draw=none, bar shift=0pt, fill=midgray] coordinates {(0.0052,1) [$0{,}005$]};
  \addplot+[draw=none, bar shift=0pt, fill=midgray] coordinates {(0.0165,2) [$0{,}016$]};
  \addplot+[draw=none, bar shift=0pt, fill=tealc]   coordinates {(0.0363,3) [$0{,}036$]};
  \addplot+[draw=none, bar shift=0pt, fill=navy!45] coordinates {(0.0063,4) [$0{,}006$]};
  \addplot+[draw=none, bar shift=0pt, fill=navy]    coordinates {(0.0077,5) [$0{,}008$]};
\end{axis}
\end{tikzpicture}

  \caption{ROC-AUC dengan CI 95\% (kiri) dan ECE (kanan) pada 28.003 respons held-out dari 492 siswa test.}
\end{figure}

\begin{table}[h]
  \centering\small
  \begin{tabular}{@{}lcccccc@{}}
    \toprule
    Model & ROC-AUC [CI 95\%] & Log-loss & Brier & Accuracy & F1 & ECE \\
    \midrule
    \rowcolor{navy!8}\textbf{P-DINA (utama)} & 0,768 [0,759--0,779] & 0,568 & 0,194 & 0,699 & 0,704 & \textbf{0,008} \\
    P-DINA prior independen & 0,761 [0,751--0,772] & 0,575 & 0,197 & 0,693 & 0,701 & 0,006 \\
    Q-logistic & 0,775 [0,766--0,786] & 0,567 & 0,193 & 0,704 & 0,726 & 0,036 \\
    Concept-score & 0,718 [0,707--0,733] & 0,611 & 0,213 & 0,658 & 0,681 & 0,016 \\
    Item-mean & 0,672 [0,661--0,682] & 0,644 & 0,226 & 0,627 & 0,665 & 0,005 \\
    \bottomrule
  \end{tabular}
  \caption{Metrik test (\texttt{artifacts/evaluation/test\_metrics.csv}).}
\end{table}

\subsection{Kalibrasi dan kualitas diagnosis}

\begin{figure}[h]
  \centering
  \setlength\figH{5.6cm}
  \begin{minipage}{0.5\linewidth}% Reliability diagram on test held-out responses; artifacts/evaluation/calibration_test.csv
\begin{tikzpicture}
\begin{axis}[
  width=\linewidth, height=\figH, xmin=0, xmax=1, ymin=0, ymax=1,
  xlabel={rata-rata prediksi P(benar)}, ylabel={proporsi benar teramati},
  grid=major, legend pos=north west, legend cell align=left,
  xtick={0,0.2,...,1}, ytick={0,0.2,...,1},
  ticklabel style={/pgf/number format/fixed, /pgf/number format/precision=1},
]
  \addplot[dashed, midgray, domain=0:1] {x};
  \addlegendentry{kalibrasi sempurna}
  \addplot[navy, thick, mark=*, mark size=1.4pt] coordinates {
    (0.0754,0.0000) (0.1721,0.1808) (0.2588,0.2608) (0.3445,0.3536) (0.4470,0.4575)
    (0.5467,0.5436) (0.6497,0.6354) (0.7476,0.7628) (0.8541,0.8502) (0.9414,0.9465)};
  \addlegendentry{P-DINA (ECE 0,008)}
  \addplot[tealc, thick, mark=square*, mark size=1.3pt] coordinates {
    (0.0780,0.1940) (0.1572,0.2362) (0.2525,0.2615) (0.3517,0.3330) (0.4499,0.3950)
    (0.5500,0.5172) (0.6492,0.6279) (0.7500,0.7755) (0.8476,0.9073) (0.9346,0.9779)};
  \addlegendentry{Q-logistic (ECE 0,036)}
\end{axis}
\end{tikzpicture}
\end{minipage}\hfill
  \begin{minipage}{0.46\linewidth}\small
    \begin{tabularx}{\linewidth}{@{}L r@{}}
      \toprule
      Ukuran & Hasil \\
      \midrule
      Konsistensi split-half (test, 880 sel) & \textbf{0,876} \\
      Akurasi atribut, simulasi (ada / tanpa bukti) & 0,957 / 0,851 \\
      Recovery parameter, simulasi $r(g)$ / $r(s)$ & 0,965 / 0,982 \\
      \bottomrule
    \end{tabularx}
  \end{minipage}
  \caption{Diagram reliabilitas (kiri) dan proksi kualitas diagnosis (kanan).}
\end{figure}

Kurva kalibrasi P-DINA hampir berimpit dengan diagonal: bila model memprediksi peluang benar 0,7, sekitar 70\% respons
memang benar. Q-logistic melenceng di kedua ujung. Konsistensi split-half 0,876 berarti dua diagnosis yang dibuat dari
dua paruh respons yang berbeda sepakat pada 88\% sel siswa--konsep. Angka simulasi adalah batas atas optimistis karena
mengasumsikan model benar.

Pada 492 siswa test (5.412 sel siswa $\times$ konsep target), diagnosis menghasilkan 727 sel \emph{mastered},
1.394 \emph{not\_mastered} (learning gap), dan 3.291 \emph{insufficient\_evidence}; 322 siswa punya setidaknya satu
gap.

\subsection{Alasan menggunakan P-DINA}

Q-logistic sedikit lebih unggul pada AUC (0,775 vs 0,768). Selisih berpasangan dari bootstrap per siswa adalah AUC
$-0{,}007$ [$-0{,}011$; $-0{,}003$] dan log-loss $+0{,}001$ [$-0{,}003$; $+0{,}007$]; artinya, dari sisi log-loss
kedua model tidak berbeda bermakna. P-DINA dipertahankan karena tiga alasan:
\begin{enumerate}
  \item probabilitasnya jauh lebih terkalibrasi (ECE 0,008 vs 0,036), padahal keputusan pendidikan memakai nilai
        probabilitas itu secara langsung;
  \item P-DINA punya konsep ``menguasai/belum menguasai'' per konsep, sedangkan Q-logistic bersifat kompensatori;
  \item parameternya bermakna langsung bagi guru (guess dan slip per soal).
\end{enumerate}
Trade-off AUC ini dilaporkan apa adanya, tidak disembunyikan. P-DINA jauh lebih baik dari concept-score dan
item-mean.

Model yang sama juga dipakai untuk asesmen adaptif (Gambar~\ref{fig:adaptif}). Pada replay 300 siswa test, 10 soal
yang dipilih secara adaptif menghasilkan label yang sama dengan diagnosis dari semua respons pada 86\% kasus,
dibanding 82\% bila soal dipilih acak.

\begin{figure}[h]
  \centering
  \setlength\figH{4.6cm}
  \begin{minipage}{0.6\linewidth}% Adaptive (mutual information) vs random item selection, replay on 300 test students;
% artifacts/evaluation/adaptive_replay_summary.csv
\begin{tikzpicture}
\begin{axis}[
  scale only axis, width=\dimexpr\linewidth-1.5cm\relax, height=\figH, xmin=4, xmax=21, ymin=0.75, ymax=0.93,
  xtick={5,10,15,20}, xlabel={panjang tes (jumlah soal)},
  ylabel={kesepakatan label vs.\ semua respons}, grid=major,
  legend pos=south east, legend cell align=left,
  yticklabel style={/pgf/number format/fixed, /pgf/number format/precision=2},
]
  \addplot[accent, thick, mark=*, mark size=1.5pt] coordinates {(5,0.789) (10,0.862) (15,0.892) (20,0.911)};
  \addlegendentry{adaptif (mutual information)}
  \addplot[midgray, thick, mark=square*, mark size=1.4pt] coordinates {(5,0.771) (10,0.824) (15,0.851) (20,0.871)};
  \addlegendentry{acak}
\end{axis}
\end{tikzpicture}
\end{minipage}
  \caption{Asesmen adaptif (mutual information) vs.\ acak (\texttt{adaptive\_replay\_summary.csv}).}
  \label{fig:adaptif}
\end{figure}

% ====================================================================================================
\section{Hyperparameter tuning}

\subsection{Seleksi model}

Grid yang diuji: prior \emph{independent} atau \emph{higher-order}, kekuatan prior Beta $\kappa\in\{2,10,40\}$, dan
class weighting ya/tidak, total 12 konfigurasi P-DINA ditambah baseline. Kriteria seleksi adalah log-loss respons
held-out pada validation.

\begin{figure}[h]
  \centering
  \setlength\figH{6.2cm}
  \begin{minipage}{0.8\linewidth}% Validation held-out log-loss per configuration; artifacts/evaluation/hp_search_validation.csv
\begin{tikzpicture}
\begin{axis}[
  xbar, scale only axis, width=\dimexpr\linewidth-3.6cm\relax, height=\figH, bar width=5.5pt, scaled x ticks=false,
  xmin=0.570, xmax=0.5848, ymin=0.4, ymax=13.6, ytick={1,...,13},
  yticklabels={
    {indep., $\kappa{=}40$, cw}, {indep., $\kappa{=}2$, cw}, {indep., $\kappa{=}10$, cw},
    {indep., $\kappa{=}40$}, {indep., $\kappa{=}10$}, {indep., $\kappa{=}2$},
    {HO, $\kappa{=}40$, cw}, {HO, $\kappa{=}10$, cw}, {HO, $\kappa{=}2$, cw},
    {Q-logistic $\sigma{=}1$}, {HO, $\kappa{=}40$}, {HO, $\kappa{=}10$}, {\textbf{HO, $\kappa{=}2$ (terpilih)}}},
  xlabel={log-loss validation (respons held-out, makin kecil makin baik)}, xmajorgrids,
  xticklabel style={/pgf/number format/fixed, /pgf/number format/precision=3},
  nodes near coords, point meta=x,
  every node near coord/.append style={font=\tiny, /pgf/number format/fixed, /pgf/number format/precision=4},
]
  \addplot+[fill=midgray, draw=none, bar shift=0pt] coordinates
    {(0.5818,1) (0.5811,2) (0.5810,3) (0.5795,4) (0.5788,5) (0.5788,6)};
  \addplot+[fill=navy!45, draw=none, bar shift=0pt] coordinates
    {(0.5745,7) (0.5742,8) (0.5742,9) (0.5726,11) (0.5721,12)};
  \addplot+[fill=tealc, draw=none, bar shift=0pt] coordinates {(0.5735,10)};
  \addplot+[fill=accent, draw=none, bar shift=0pt] coordinates {(0.5721,13)};
\end{axis}
\end{tikzpicture}
\end{minipage}
  \caption{Log-loss validation per konfigurasi (\texttt{hp\_search\_validation.csv}). HO = higher-order,
  cw = class weighting.}
\end{figure}

Temuannya: prior higher-order konsisten lebih baik (0,572 vs 0,579); kekuatan prior $\kappa$ hampir tidak berpengaruh
karena data per soal sudah cukup banyak; class weighting memperburuk log-loss dan kalibrasi di semua konfigurasi
sehingga tidak dipakai. Konfigurasi terpilih: \textbf{higher-order, $\kappa=2$, tanpa class weighting}.

\subsection{Kebijakan diagnosis}

P(mastery) diubah menjadi label dengan tiga angka yang juga dipilih dari validation.

\begin{figure}[h]
  \centering
  \setlength\figH{4.6cm}
  \begin{minipage}{0.48\linewidth}% Threshold tau selection on validation; artifacts/evaluation/threshold_selection_validation.csv
\begin{tikzpicture}
\begin{axis}[
  width=\linewidth, height=\figH, xmin=0.28, xmax=0.72, ymin=0.5850, ymax=0.5910,
  xlabel={threshold $\tau$}, ylabel={log-loss held-out (label keras)},
  grid=major, xtick={0.3,0.4,...,0.7},
  xticklabel style={/pgf/number format/fixed, /pgf/number format/precision=2},
  yticklabel style={/pgf/number format/fixed, /pgf/number format/precision=4},
  scaled y ticks=false,
]
  \fill[accent!12] (axis cs:0.425,0.5850) rectangle (axis cs:0.72,0.5910);
  \node[font=\tiny, text=accent!80!black, anchor=north] at (axis cs:0.575,0.5908) {dalam 1 SE dari yang terbaik};
  \addplot[navy, thick, mark=*, mark size=1.4pt] coordinates {
    (0.30,0.59024) (0.35,0.58909) (0.40,0.58826) (0.45,0.58661) (0.50,0.58607)
    (0.55,0.58612) (0.60,0.58607) (0.65,0.58583) (0.70,0.58609)};
  \addplot[only marks, mark=*, mark size=3pt, accent] coordinates {(0.50,0.58607)};
  \node[font=\tiny, text=accent!80!black, anchor=north east] at (axis cs:0.49,0.58600) {$\tau{=}0{,}5$ dipilih};
  \node[font=\tiny, anchor=north] at (axis cs:0.65,0.58575) {minimum mentah 0,65};
\end{axis}
\end{tikzpicture}
\end{minipage}\hfill
  \begin{minipage}{0.48\linewidth}% Minimum evidence n_min from split-half classification consistency on validation (README_pdina.md, section 6)
\begin{tikzpicture}
\begin{axis}[
  ybar, width=\linewidth, height=\figH, bar width=9pt,
  ymin=0.78, ymax=0.90, xtick={1,2,3,4,5,6,8}, xmin=0.3, xmax=8.7,
  xlabel={$n_{\min}$ (respons per konsep, per paruh)}, ylabel={konsistensi split-half},
  ymajorgrids, nodes near coords, point meta=y,
  every node near coord/.append style={font=\tiny, /pgf/number format/fixed, /pgf/number format/precision=3},
  yticklabel style={/pgf/number format/fixed, /pgf/number format/precision=2},
]
  \addplot+[fill=navy!45, draw=none, bar shift=0pt] coordinates {(1,0.837) (2,0.825) (3,0.820) (4,0.839) (6,0.879) (8,0.881)};
  \addplot+[fill=accent, draw=none, bar shift=0pt] coordinates {(5,0.857)};
  \draw[gapred, dashed, thick] (axis cs:0.3,0.85) -- (axis cs:8.7,0.85);
  \node[font=\tiny, text=gapred, anchor=south] at (axis cs:7,0.85) {target 0,85};
\end{axis}
\end{tikzpicture}
\end{minipage}
  \caption{Pemilihan threshold $\tau$ (kiri) dan bukti minimum $n_{\min}$ (kanan) pada validation.}
\end{figure}

\begin{itemize}
  \item $\boldsymbol{\tau=0{,}5}$. Nilai terbaik mentah adalah 0,65, tetapi semua $\tau$ di rentang 0,45--0,70 berada
        dalam satu standard error dari yang terbaik. Aturan \emph{one-standard-error} memilih nilai netral 0,5.
  \item $\boldsymbol{n_{\min}=5}$. Respons tiap siswa dibagi dua secara acak, diagnosis dibuat terpisah, lalu
        kesepakatannya dihitung. Lima adalah nilai terkecil yang mencapai konsistensi minimal 0,85.
  \item $\boldsymbol{\delta=0}$. Target konsistensi sudah tercapai tanpa pita ketidakpastian, sehingga label
        \emph{uncertain} jarang muncul.
\end{itemize}

\begin{table}[h]
  \centering\small
  \begin{tabular}{@{}ll@{}}
    \toprule
    Label & Aturan \\
    \midrule
    \emph{insufficient\_evidence} & respons pada konsep $<n_{\min}$ (butuh soal tambahan, bukan gap) \\
    \emph{mastered} & $\pi\ge\tau+\delta$ \\
    \emph{not\_mastered} (= learning gap) & $\pi<\tau-\delta$ \\
    \emph{uncertain} & $|\pi-\tau|<\delta$ \\
    \bottomrule
  \end{tabular}
  \caption{Aturan label diagnosis.}
\end{table}

% ====================================================================================================
\section{Penyimpanan model}

Hasil training disimpan sebagai checkpoint di \texttt{artifacts/model/} dengan total sekitar 1,1 MB. Tidak ada bobot
jaringan saraf atau framework deep learning.

\begin{table}[h]
  \centering\small
  \begin{tabular}{@{}>{\ttfamily}l P{6.2cm} r@{}}
    \toprule
    \normalfont File & Isi & Ukuran \\
    \midrule
    pdina\_params.npz & Q-matrix, $g$, $s$, $\lambda$, $a$, $b$, ID soal dan konsep & 459 KB \\
    pdina\_config.json & hyperparameter, jumlah iterasi EM, ELBO akhir & 1 KB \\
    diagnosis\_policy.json & $\tau$, $\delta$, $n_{\min}$, konsep target, dan asal-usul pemilihannya & 1 KB \\
    item\_parameters.csv & guess, slip, daya beda, dan kelayakan per soal & 202 KB \\
    item\_bank.csv, q\_matrix.csv & item bank dan Q-matrix yang dipakai & 331 KB \\
    concept\_metadata.csv, \dots & nama konsep, cakupan target, prasyarat, urutan kurikulum & 23 KB \\
    intervention\_map\_resolved.csv & peta intervensi (placeholder) & 8 KB \\
    model\_card.json & ringkasan run, data, dan metrik & 17 KB \\
    \bottomrule
  \end{tabular}
  \caption{Isi checkpoint.}
\end{table}

Saat dipakai, \texttt{CDMService.load()} memuat checkpoint dan memeriksa konsistensi ID soal, urutan konsep, dan
konsep target. Fungsi \texttt{diagnose()} hanya menjalankan E-step dengan parameter beku. Aplikasi Streamlit tidak
membaca dataset jutaan baris dan tidak melatih ulang.

Training ulang diperlukan bila ada soal baru (parameter $g$ dan $s$ harus dikalibrasi dari minimal sekitar 30 respons
per soal), bila Q-matrix atau daftar konsep berubah, atau bila model dipakai untuk populasi lain seperti kelas SMA
penelitian. Mengubah $\tau$, $\delta$, $n_{\min}$, atau peta intervensi tidak memerlukan training ulang.

% ====================================================================================================
\section{Pemilihan LLM dan spesifikasi}

Versi awal sistem memakai \emph{small language model} lokal, Qwen3-4B-Instruct, yang dijalankan dengan Ollama di
laptop (GPU RTX 2050 4 GB). Model itu menulis angka dengan benar, tetapi satu laporan guru membutuhkan sekitar 80--105
detik, memenuhi VRAM, dan membuat baterai cepat habis. Karena tugas LLM di sistem ini hanya menulis ulang fakta yang
sudah dihitung, dan kebenaran angka dijaga validator, model kecil yang cepat sudah cukup. Pilihannya jatuh pada
\textbf{Gemini 3.1 Flash-Lite} lewat Gemini API resmi.

\begin{table}[h]
  \centering\small
  \begin{tabularx}{\linewidth}{@{}l L L@{}}
    \toprule
    & SLM lokal (sebelumnya) & \textbf{Gemini 3.1 Flash-Lite} \\
    \midrule
    Model & Qwen3-4B-Instruct (Ollama) & \texttt{gemini-3.1-flash-lite} \\
    Tempat jalan & laptop, RTX 2050 4 GB & server Google (Gemini API) \\
    Waktu per laporan guru & sekitar 80--105 detik & 5,97 detik (30 siswa demo) \\
    Validator & lolos pada percobaan pertama & lolos pada percobaan pertama \\
    Beban laptop & VRAM penuh, baterai boros & hanya klien HTTP \\
    Kebutuhan & unduh model 2,5 GB, server lokal & API key dan internet \\
    \bottomrule
  \end{tabularx}
  \caption{Perbandingan yang sudah diukur di proyek ini.}
\end{table}

Spesifikasi yang dipakai: akses lewat Gemini API dengan key dari Google AI Studio, SDK resmi \texttt{google-genai}
versi 2.33.0, model \texttt{gemini-3.1-flash-lite} (versi stabil; versi preview juga tersedia di akun), mode teks ke
teks tanpa streaming, thinking level \texttt{low}, dan satu request per laporan (ditambah repair bila gagal validasi).
Pada laporan guru untuk 30 siswa demo EEDI, pemakaian terukur adalah 3.246 token input, 618 token output, dan 119 token
thinking, dengan latensi 5,97 detik dalam satu percobaan.

\begin{peringatan}
  Fakta yang dikirim sudah dipseudonimkan dan diagregasi, tetapi tetap dikirim ke server Google. Untuk data kelas
  nyata, gunakan project dengan billing aktif dan periksa ketentuan penggunaan data serta kebijakan sekolah dan etika
  penelitian.
\end{peringatan}

% ====================================================================================================
\section{Konfigurasi model}

Semua pengaturan LLM ada di \texttt{configs/llm.json}.

\begin{table}[h]
  \centering\small
  \begin{tabular}{@{}>{\ttfamily}l >{\ttfamily}l P{6.6cm}@{}}
    \toprule
    \normalfont Kunci & \normalfont Nilai & Alasan \\
    \midrule
    model & gemini-3.1-flash-lite & versi stabil Flash-Lite \\
    api\_key\_env & GEMINI\_API\_KEY & key hanya dibaca dari environment \\
    env\_file & .env & cadangan bila environment proses kosong; file ini di-\texttt{.gitignore} \\
    temperature & 1.0 & rekomendasi Google untuk keluarga Gemini 3 \\
    top\_p & null & memakai default API \\
    thinking\_level & low & menekan latensi dan biaya \\
    max\_tokens & 8192 & batas output, termasuk token thinking \\
    timeout\_sec & 120 & batas waktu per request \\
    max\_retries & 2 & jumlah repair setelah gagal validasi \\
    prompt\_version & narrative-gemini-v1 & ikut menentukan kunci cache \\
    max\_class\_concepts & 11 & batas ukuran fakta; gap dan rekomendasi masing-masing 5 \\
    max\_examples & 1 & contoh few-shot per request \\
    cache\_ttl\_hours & 168 & laporan yang sama diambil dari cache dalam kurang dari 1 detik \\
    \bottomrule
  \end{tabular}
  \caption{Konfigurasi LLM.}
\end{table}

Sistem juga dirancang tahan gangguan: HTTP 429 dan 5xx di-retry dengan backoff 2 dan 8 detik; tiga kegagalan
berturut-turut membuka \emph{circuit breaker} selama 60 detik; output yang terpotong (\texttt{MAX\_TOKENS}) atau kosong
ditolak; tanpa key atau internet, laporan template langsung ditampilkan tanpa error. Audit log mencatat waktu, model,
hash input, latensi, token, status validasi, dan fallback, tetapi tidak menyimpan teks laporan, ID siswa, maupun
API key.

% ====================================================================================================
\section{System prompt dan pengaman}

\subsection{System prompt}

Pesan yang dikirim ke Gemini terdiri atas tiga bagian: \texttt{system\_instruction} berisi aturan; satu contoh
few-shot (giliran user lalu model) untuk audiens yang sama; dan pesan user berisi ``FAKTA DIAGNOSIS (JSON)'' beserta
instruksi menulis. Berikut teks system prompt yang persis dipakai untuk audiens guru:

\lstinputlisting{data/system_prompt_teacher.txt}

Untuk audiens siswa, prompt meminta sapaan ``kamu'', melarang menyebut data kelas atau siswa lain, dan memakai heading
Ringkasan, Yang sudah dikuasai, Yang perlu dilatih, Langkah belajar, dan Catatan.

\subsection{Pengaman di sekitar LLM}

Prompt hanya \emph{meminta}; kepatuhan \emph{diperiksa} oleh validator (Gambar~\ref{fig:guard}).

\begin{figure}[h]
  \centering
  \resizebox{0.95\linewidth}{!}{% LLM guardrail: validator, repair, deterministic fallback.
\begin{tikzpicture}[node distance=5mm and 8mm]
  \node[boxsoft, text width=1.9cm] (f) {Fakta JSON\\{\tiny dari P-DINA}};
  \node[boxsoft, right=of f, text width=2.2cm] (p) {Prompt\\{\tiny aturan + heading + 1 contoh}};
  \node[boxacc, right=of p, text width=1.9cm] (g) {\textbf{Gemini}\\{\tiny 1 request}};
  \node[boxnavy, right=of g, text width=1.9cm] (v) {\textbf{Validator}\\{\tiny hard fail}};
  \node[box, draw=okgreen, right=of v, text width=2.0cm] (ok) {\textbf{Tampilkan}\\{\tiny + cache disk}};
  \node[boxsoft, below=6mm of g, text width=2.4cm] (r) {Repair\\{\tiny catat pelanggaran, maks.\ 2$\times$}};
  \node[box, draw=gapred, below=6mm of ok, text width=2.0cm] (t) {\textbf{Template}\\{\tiny deterministik}};
  \draw[arr] (f) -- (p);
  \draw[arr] (p) -- (g);
  \draw[arr] (g) -- (v);
  \draw[arr] (v) -- node[lbl, above]{lolos} (ok);
  \draw[arr] (v) |- node[lbl, pos=0.25, right]{gagal} (r);
  \draw[arr] (r) -- (g);
  \draw[arr] (r.east) ++(0,-0.15) -- node[lbl, below, pos=0.6]{masih gagal / API mati} ($(t.west)+(0,-0.15)$);
\end{tikzpicture}
}
  \caption{Validator, repair, dan fallback template.}
  \label{fig:guard}
\end{figure}

Validator menolak teks bila: ada angka yang tidak berasal dari fakta JSON, atau angka pada baris yang menyebut sebuah
konsep ternyata milik konsep lain; heading wajib tidak lengkap; ada HTML atau daftar bernomor; rekomendasi placeholder
tidak disebut ``perlu validasi guru''; muncul klaim terlarang (\emph{terbukti, dijamin, efektif meningkatkan, siswa
malas}, dan sebagainya); ada konsep di luar fakta; atau panjangnya di luar 120--900 kata. Bila gagal, Gemini diminta
memperbaiki dengan daftar pelanggarannya (maksimal dua kali). Bila tetap gagal atau API tidak tersedia, sistem
menampilkan laporan template deterministik yang selalu lolos validator. Dengan demikian pengguna tidak pernah melihat
teks LLM yang belum divalidasi.

Dari sisi privasi, fakta yang dikirim tidak memuat matriks respons, nama, sekolah, atau \texttt{UserId} asli. Siswa
disebut \texttt{Siswa-<hash>}, dan laporan guru hanya berisi agregat kelas.

% ====================================================================================================
\section{Output sistem AI secara keseluruhan}

\subsection{Diagnosis P-DINA}

Contoh berikut memakai 30 siswa pertama dari data demo (siswa test EEDI) dengan checkpoint EEDI: 8.653 respons pada
11 konsep target.

\begin{figure}[h]
  \centering
  \setlength\figH{5.6cm}
  \begin{minipage}{0.75\linewidth}% Class summary of 30 EEDI demo students (facts sent to the LLM); slides/data/example_teacher_report.json
\begin{tikzpicture}
\begin{axis}[
  xbar stacked, scale only axis, width=\dimexpr\linewidth-4.0cm\relax, height=\figH, bar width=6pt,
  xmin=0, xmax=30, ytick={1,...,11}, enlarge y limits=0.06,
  yticklabels={Gradient Between Two Co-ordinates, Naming Co-ordinates in 2D, Factorising into a Single Bracket,
               Writing Expressions, Simplifying Expressions (Like Terms), Horizontal and Vertical Lines,
               Substitution into Formula, Solving Linear Inequalities, Midpoint Between Two Co-ordinates,
               Linear Sequences (nth term), Linear Equations},
  xlabel={jumlah siswa (dari 30)}, xmajorgrids,
  legend style={at={(0.5,1.03)}, anchor=south, legend columns=3, font=\tiny, /tikz/every even column/.append style={column sep=4pt}},
  y tick label style={font=\tiny},
]
  \addplot+[fill=gapred, draw=none] coordinates {(0,1) (8,2) (9,3) (12,4) (13,5) (5,6) (12,7) (4,8) (4,9) (6,10) (9,11)};
  \addlegendentry{learning gap (\emph{not\_mastered})}
  \addplot+[fill=okgreen, draw=none] coordinates {(0,1) (8,2) (7,3) (8,4) (8,5) (3,6) (7,7) (2,8) (2,9) (3,10) (4,11)};
  \addlegendentry{\emph{mastered}}
  \addplot+[fill=insuf, draw=none] coordinates {(30,1) (14,2) (14,3) (10,4) (9,5) (22,6) (11,7) (24,8) (24,9) (21,10) (17,11)};
  \addlegendentry{bukti kurang}
\end{axis}
\end{tikzpicture}
\end{minipage}
  \caption{Ringkasan kelas: jumlah siswa per label untuk tiap konsep target.}
\end{figure}

\begin{table}[h]
  \centering\footnotesize\setlength\tabcolsep{5pt}
  \begin{tabular}{@{}lrrrrl@{}}
    \toprule
    Konsep & Bukti & Benar & P(mastery) & SD & Label \\
    \midrule
    Linear Sequences (nth term) & 11 & 1 & 0,00 & 0,00 & learning gap \\
    Simplifying Expressions by Collecting Like Terms & 19 & 8 & 0,00 & 0,00 & learning gap \\
    Substitution into Formula & 22 & 8 & 0,00 & 0,01 & learning gap \\
    Factorising into a Single Bracket & 8 & 3 & 0,00 & 0,02 & learning gap \\
    Linear Equations & 7 & 3 & 0,00 & 0,03 & learning gap \\
    Solving Linear Inequalities & 8 & 4 & 0,03 & 0,18 & learning gap \\
    Naming Co-ordinates in 2D & 11 & 4 & 0,04 & 0,20 & learning gap \\
    Writing Expressions & 25 & 13 & 0,34 & 0,47 & learning gap \\
    Horizontal and Vertical Lines & 4 & 0 & 0,00 & 0,01 & bukti belum cukup \\
    Midpoint Between Two Co-ordinates & 4 & 1 & 0,00 & 0,02 & bukti belum cukup \\
    Gradient Between Two Co-ordinates & 1 & 1 & 0,47 & 0,50 & bukti belum cukup \\
    \bottomrule
  \end{tabular}
  \caption{Profil satu siswa (\texttt{Siswa-e7f6c011}). SD = $\sqrt{\pi(1-\pi)}$, ukuran ketidakpastian, bukan
  interval kepercayaan.}
\end{table}

Perhatikan tiga konsep terakhir: walaupun probabilitasnya rendah, konsep dengan bukti kurang dari lima respons tidak
dicap learning gap. Untuk setiap gap, aturan R1--R5 menghasilkan tindakan (remediasi, asesmen ulang, atau tambah
bukti), level intervensi (seluruh kelas, kelompok kecil, atau individual), aktivitas, dan soal latihan dari item bank.
Isi aktivitas masih placeholder sampai divalidasi guru.

\subsection{Laporan naratif Gemini}

Fakta kelas di atas dikirim ke \texttt{gemini-3.1-flash-lite}. Laporan berikut lolos validator pada percobaan pertama
(5,97 detik; 3.246 token input, 618 output, 119 thinking). Kutipan apa adanya (dipotong):

\begin{tcolorbox}[colback=white, colframe=accent, boxrule=0.5pt, arc=2pt, fontupper=\small,
                  title={\small Laporan guru -- dihasilkan \texttt{gemini-3.1-flash-lite}}, coltitle=white, colbacktitle=accent]
  \textbf{Ringkasan kelas}\\
  Data mencakup 30 siswa yang diuji pada 11 konsep target dengan total 8653 respons. Estimasi model mendeteksi 18
  siswa memiliki learning gap berdasarkan ambang batas diagnostik yang ditetapkan.

  \textbf{Temuan utama}
  \begin{itemize}
    \item \textbf{Simplifying Expressions by Collecting Like Terms} memiliki 13 siswa yang belum menguasai, dengan
          estimasi probabilitas penguasaan 35\%.
    \item \textbf{Substitution into Formula} menunjukkan 12 siswa belum menguasai, dengan estimasi probabilitas
          penguasaan 30\%.
    \item \textbf{Writing Expressions} tercatat 12 siswa belum menguasai, dengan estimasi probabilitas penguasaan 42\%.
    \item \textbf{Linear Equations} memiliki 9 siswa dengan learning gap, dengan estimasi probabilitas penguasaan 30\%.
    \item \textbf{Gradient Between Two Co-ordinates} mencatat 30 siswa dengan bukti tidak cukup, yang berarti
          memerlukan soal tambahan untuk mendapatkan diagnosis akurat.
  \end{itemize}

  \textbf{Rekomendasi tindakan}
  \begin{itemize}
    \item \textbf{Simplifying Expressions by Collecting Like Terms} (seluruh kelas): Penjelasan ulang dengan contoh
          terbimbing, lalu latihan bertahap dengan umpan balik untuk 13 siswa. Rekomendasi ini adalah placeholder dan
          perlu validasi guru. [\dots]
  \end{itemize}

  \textbf{Keterbatasan dan catatan}
  \begin{itemize}
    \item Seluruh data penguasaan adalah estimasi model dari respons siswa, bukan hasil terverifikasi.
    \item Status bukti tidak cukup (insufficient evidence) menunjukkan perlunya tambahan data, bukan sebagai learning gap.
  \end{itemize}
\end{tcolorbox}

Semua angka dalam laporan (30, 11, 8653, 18, 13, 35\%, dan seterusnya) berasal dari fakta P-DINA; persentase seperti
35\% adalah rata-rata P(mastery) konsep itu (0,35) yang ditulis sebagai persen. Validator memeriksa angka dan format,
tetapi tidak menilai ketepatan pedagogis, sehingga laporan tetap diperlakukan sebagai draf yang ditinjau guru.

% ====================================================================================================
\section{Kesimpulan}

\begin{enumerate}
  \item P-DINA dengan prior higher-order dan variational EM dilatih pada 1,1 juta respons train EEDI dalam sekitar
        30 detik di CPU, dan menghasilkan probabilitas penguasaan per siswa dan konsep yang terkalibrasi (ECE 0,008).
  \item Akurasi prediksi responsnya setara dengan baseline terkuat (AUC 0,768 vs 0,775; log-loss tidak berbeda
        bermakna), dengan parameter yang dapat dijelaskan.
  \item Kebijakan diagnosis dipilih dari data validation ($\tau=0{,}5$, $n_{\min}=5$, $\delta=0$), dengan konsistensi
        split-half test 0,876.
  \item Gemini 3.1 Flash-Lite hanya mengomunikasikan hasil diagnosis: sekitar 6 detik per laporan dan lolos validator
        pada percobaan pertama. Tanpa API key, sistem tetap berjalan dengan laporan template.
  \item Seluruh sistem berjalan end-to-end di laptop berbasis CPU dan dikemas dalam aplikasi Streamlit.
\end{enumerate}

% ====================================================================================================
\section{Keterbatasan dan langkah berikutnya}

\subsection{Keterbatasan}
\begin{itemize}
  \item EEDI berisi siswa SMP di Inggris. Ke-11 konsep target hanya prasyarat aljabar SMA dan tidak ada kalkulus.
  \item Tidak ada ground truth mastery; kualitas diagnosis hanya diukur lewat proksi.
  \item Q-matrix diasumsikan benar, padahal 5 dari 22 soal yang diaudit hanya cocok sebagian.
  \item Daya beda soal sedang: 3.291 dari 5.412 sel test masih berlabel \emph{insufficient\_evidence}, dan konsep
        \emph{Gradient Between Two Co-ordinates} hanya punya 2 soal sehingga tidak dapat didiagnosis.
  \item Rekomendasi intervensi masih placeholder dan tidak ada klaim efektivitas.
  \item Laporan LLM baru diuji pada beberapa laporan nyata, bergantung pada internet dan kuota, serta mengirim fakta
        (pseudonim) ke Google.
\end{itemize}

\subsection{Langkah berikutnya}
\begin{enumerate}
  \item Menyusun bank soal SMA (aljabar dan kalkulus) dengan Q-matrix yang ditinjau guru.
  \item Mengumpulkan respons kelas nyata dan melatih ulang model (minimal sekitar 30 respons per soal).
  \item Menjalankan pre-test dan post-test dengan kelompok pembanding.
  \item Meminta guru memvalidasi peta intervensi.
  \item Mengevaluasi laporan LLM secara sistematis dengan \texttt{scripts/eval\_narrative.py} dan menilainya bersama
        guru.
\end{enumerate}

\vspace{1em}
\noindent Terima kasih.

% ====================================================================================================
\appendix
\section{Referensi}
\begin{itemize}
  \item Junker, B. W., \& Sijtsma, K. (2001). Cognitive assessment models with few assumptions, and connections with
        nonparametric item response theory. \emph{Applied Psychological Measurement}, 25(3), 258--272.
  \item de la Torre, J., \& Douglas, J. A. (2004). Higher-order latent trait models for cognitive diagnosis.
        \emph{Psychometrika}, 69(3), 333--353.
  \item Blei, D. M., Kucukelbir, A., \& McAuliffe, J. D. (2017). Variational inference: A review for statisticians.
        \emph{Journal of the American Statistical Association}, 112(518), 859--877.
  \item Wang, Z., et al. (2020). Instructions and guide for diagnostic questions: The NeurIPS 2020 Education
        Challenge. arXiv:2007.12061.
\end{itemize}

\end{document}
